\section{Discussion}
 Prior work has shown that a close variant of the \newrecipe recipe can converge on small models over short token horizons without keeping any layers in high precision \citep{panferov2026quartetii}. However, at scale these NVFP4 recipes still require high-precision layers, suggesting a representation failure that the error-reduction techniques studied in this work cannot address. This is shown most directly on the 8B parameter model trained to 300B tokens in Figure \ref{fig:last_15_pct_bf16}, and further supported by Figure \ref{fig:backward_unbiased}, where all recipes, regardless of the backward-tensor technique used, show a loss gap that grows throughout training.

 Figure \ref{fig:recipe_stacking} shows that the relative ranking of recipes does not stabilize until well past 150B tokens of training. Notably, weaker recipes often initially appear better than the strongest ones. These early changes in ordering can be attributed to differences in the randomly sampled sign vectors used by the Random Hadamard Transforms, noise in training, or convergence properties of individual recipes. Only after enough tokens have been seen do the rankings settle. This view suggests that short NVFP4 pretraining ablations are useful for detecting catastrophic instability, but they are unfortunately not sufficient for deciding whether a new training algorithm will preserve accuracy at scale. For FP4 pretraining methods, where the intended benefit is to make long training runs cheaper, this late-emerging failure mode is especially important.

 The separation between techniques applied to backward tensors is more visible in a less stable training regime. Figure \ref{fig:backward_unbiased} shows a slight separation when all layers are quantized, but indistinguishable curves once the last 15\% of the network is kept in BF16. For our full training runs, we adopted a recipe that addresses bias in the gradients and reduces quantization error in the subsequent backward tensors, applying targeted interventions to each tensor type. A pragmatic alternative entails applying MS-EDEN to gradients without additional techniques on X.T and W.T, likely without significant cost to model quality in the stable regime, as Figure \ref{fig:backward_unbiased} suggests.

The \newrecipe recipe incrementally improves upon the \nvrecipe recipe, remains robust across the two model sizes, datasets, and architectures explored in this work, and partially closes the training-loss gap to BF16 at both scales.

